\chapter{FUTURE ENHANCEMENT}

\hspace{0.5cm} The AI-Powered Code Document Generator sets a robust foundation for automated documentation, demonstrating how static code analysis and large language models can work together to produce accurate, structured and context-rich documentation for multiple programming languages. This chapter explores several opportunities for enhancing and refining the system to advance it toward a more comprehensive, production-ready platform.

\section{Git Integration}
\hspace{0.5cm} One of the way to improve the system would be to integrate it more closely with Git-based platforms like GitHub and GitLab. Currently, documentation is generated when a user provides a repository URL or manually uploads files. Enhanced Git integration could allow the documentation to update automatically with each commit.



\section{Multi-Language Expansion}
\hspace{0.5cm} Although the system currently supports fourteen programming languages, future upgrades could add support for languages like Swift, Scala, Dart and Haskell. Expanding language compatibility would make the system applicable to a wider range of development projects beyond its current scope.


\section{Interactive Q\&A}
\hspace{0.5cm} A highly valuable future feature could be the addition of an intelligent Question\&Answer system (Q\&A ), allowing developers to interact with the codebase through a natural language interface. Instead of merely reviewing the automatically generated documentation, developers could directly ask questions about function call chains, recent changes or other specific aspects of the system.


\newline
\hspace{0.5cm}The future improvements proposed in this chapter represent a natural progression of the current system. Git integration would embed documentation generation into the ongoing development workflow. Expanding multi-language support would broaden the system’s applicability across diverse programming environments. Adding an interactive Q\&A feature would transform the system from a passive documentation generator into an active tool for code understanding. Collectively, these enhancements would result in a highly efficient and practical automated documentation platform.